\documentclass[11pt]{article}
\textwidth 15.9cm
\textheight 23. cm
\addtolength{\oddsidemargin}{-11.5mm}
\addtolength{\evensidemargin}{-11.5mm}
\addtolength{\topmargin}{-19.0mm}

\usepackage[ansinew]{inputenc}
\usepackage{amsmath,amssymb}
\usepackage{bbm}
\usepackage{graphicx}
\usepackage{hyperref}

\include{GempicLatexMacros}

\renewcommand{\part}{\texttt{part}}
\newcommand{\cell}{\texttt{cell}}
\newcommand{\node}{\texttt{node}}
\newcommand{\prim}{\texttt{prim}}

\begin{document}

\section{Particle Structure}\label{sec:ParticleStructure}
Particles in GEMPICX are divided into vectors of shared pointers to species, or groups, based on the AMReX particle structure (see \href{https://amrex-codes.github.io/amrex/docs_html/Particle.html}{the AMReX particle documentation} for more details).
Each group contains the properties
\begin{itemize}
    \item charge
    \item mass
    \item name
    \item \lstinline{vDim}: The number of velocity dimensions
    \item \lstinline{nData}: The number of additional real data members (typically just weight)
\end{itemize}
Individual particles then have further properties
\begin{itemize}
    \item \lstinline{id} and \lstinline{cpu}: Constant AMReX-generated identifiers, the combination of which is unique to each particle
    \item \lstinline{pos}: arrays of particle positions
    \item velocities
    \item weight
    \item possibly others, depending on the simulation
\end{itemize}
Particles are accessed and manipulated through the \lstinline{amrex::ParIter} construct, where a special for-loop such as, \emph{e.g.},
\begin{lstlisting}
for (amrex::ParIter<0, 0, vDim + nData, 0> pti(*particleSpecies, 0); pti.isValid(); ++pti)
\end{lstlisting}
creates accessors for particle data and distributes work across the different processes available to the simulation. Pointers to particle data are then obtained via this particle iterator.

In GEMPICX, we currently have positions as arrays of structs and velocities and other extra data as struct of arrays.
We do not currently provide the ability to have extra integer components.
Since particle all data is stored on the GPU, the procedure is then first to obtain the pointers, launch a \lstinline{ParallelFor} loop, and access the data inside, \emph{e.g.}
\begin{lstlisting}
for (amrex::ParIter<0, 0, vDim + nData, 0> pti(*particleSpecies, 0); pti.isValid(); ++pti)
{
    // Getting particle data pointers
    const long np = pti.numParticles();
    const auto &particles = pti.GetArrayOfStructs()().data();
    auto *const velx = pti.GetStructOfArrays().GetRealData(0).data();
    auto *const vely = pti.GetStructOfArrays().GetRealData(1).data();
    auto *const velz = pti.GetStructOfArrays().GetRealData(2).data();
    auto *const weight = pti.GetStructOfArrays().GetRealData(vDim).data();

    // other possible prepwork
    ...

    // launch GPU loop over all particles in current tile
    amrex::ParallelFor(np, [=] AMREX_GPU_DEVICE(long pp)
    {
        // do stuff with particles
        particles[pp].pos(xDir) = weight[pp] + 0.5 * velx[pp];
    });
}
\end{lstlisting}


\section{Particle Samplers}
At the moment, particles can be initiated with pseudo-random, quasi-random, and rejection samplers outside of tests.
They are initialized using the \lstinline{Gempic::Particle::init_particles} function.
The different options are selected in the input file (see \ref{sec:ParticleInput}), the details of which follow below.

All the direct methods use random values to
\begin{itemize}
    \item initialize positions from a uniform distribution on the sampling region.
    \item initialize velocities according to one or more Gaussian distributions given in the input file, each specified by mean velocity and the standard deviation of velocity either as a constant or a function in all directions. Initialized by random values through the Box-Muller method.
\end{itemize}
The particle weight is proportional to the mandatory input function in the input file evaluated at the particle position. The particles are spatially uniform and the density profile is represented only in the weights.

The seed of the random generators can be specified in the input file.

\begin{note}
    All samplers can be used on CPU and GPU. \ref{sec:sampler_pseudorandomGPU} is currently the only one running natively and in parallel on device. All others sample on host to a buffer that gets flushed to the device memory. This approach is of coarse less efficient for GPU simulations but allows for easier sampler development that work out of the box on different architectures.
\end{note}

\subsection{Sampler: PseudoRandom}
Generates uniformly distributed random values for position and velocity distributions for all particles cell by cell.

\subsection{Sampler: PseudoRandomBox}
Same as \verb|PseudoRandom| but draws unformly distributed position values on boxes instead of individual cells. There is no guarantee that \verb|nPartPerCell| particles end up in each cell.

\subsection{Sampler: PseudoRandomGPU}
\label{sec:sampler_pseudorandomGPU}
Equivalent to \verb|PseudoRandom| but the sampling is done fully in parallel on device using \verb|amrex::ParallelForRNG|. Also runs in CPU but does not have any advantages in that scenario.

\subsection{Sampler: Sobol}
Positions and velocities are generated based on quasirandom numbers from the Sobol sequence. The sampling is done on the box level and each box uses a random starting index in the sequence determined by the random seed.

\subsection{Sampler: Rejection}
Sampling arbitrary target distribution functions by accepting/rejecting samples of a proposal distribution depending on the ratio of (probability) densities. The proposal distribution is always a Maxwellian that can be configured in the input file and sampled directly. The target distribution is specified by a string \verb|Particle.species.rejectionDistribution| that in addition to the spatial variables \verb|x, y, z| can contain the velocity components \verb|vx, vy, vz|. An example is the standard univariate Maxwellian

\begin{lstlisting}[language=pkgconfig]
    rejectionDistribution = "1 / sqrt(2*3.14159265)^3 * exp(-(vx^2 + vy^2 + vz^2)/(2))"
\end{lstlisting}

but functions can be more involved like a spatial temperature profile at the gyrocenters of particles in a constant magnetic field

\begin{lstlisting}[language=pkgconfig]
    rejectionDistribution = "1 / sqrt(2*3.14159265*(1+0.5*cos(kvarx*(x-vy))))^3 * exp(-(vx^2 + vy^2 + vz^2)/(2*(1+0.5* cos(kvarx*(x-vy)))))"
\end{lstlisting}

To ensure that $target <= proposal$, the proposal is scaled by \verb|Particle.species.rejectionScaling|. To avoid infinite loops, the sampling aborts if \verb|Particle.species.nMaxRejections| samples get rejected in a row. After sampling, some diagnostic information is returned. In addition to the acceptance rate it contains the minimal scaling factor $s$ such that $target <= s \cdot proposal$:

\begin{lstlisting}
    Rejection sampling: Acceptance rate                = 0.03997958542
    Rejection sampling: Samples with target > proposal = 0
    Rejection sampling: Minimum scaling necessary      = 23.21383126
\end{lstlisting}

When using parallelization, each process samples globally followed by a particle redistribution.

\begin{note}
    The spatial density of the proposal distribution should be set to \verb|density = 1| to prevent unwanted effects and obtain the expected spatial density of the specified distribution function.

    There are currently the following restrictions on the proposal distribution
    \begin{itemize}
        \item Only a single Gaussian
    \end{itemize}
\end{note}

\subsection{Sampler: RejectionSobol}
Same as \verb|Rejection| but using quasi-random values from the Sobol sequence.

\end{document}